% Kamp III.2, Lemma 5: existential composition through a shared midpoint.

\subsubsection{Composition through an interior point}
\label{sec:p09-kamp-iii2-lemma5}

Kamp's Lemma~5 combines elementary conditions on two intervals meeting at
one quantified point. The point must lie strictly between the outer
endpoints and satisfy its own elementary conditions. The source uses
Lemma~4 for the conjunction on each side, distributes over their
elementary disjuncts, and concatenates each pair of finite patterns
\cite[ch.~III, lemma~5, pp.~78--79]{Kamp1968}.

\begin{hbdefinition}{def:p09-kamp-midpoint-word}{Concatenating two words at an elementary midpoint}
Use the elementary labelled words and formula recursion
\(\mathcal B_p\) of \cref{def:p09-kamp-labelled-intervals}.
Fix a common free parameter \(w\), a label
\(\eta\in\operatorname{E}(w)\), and words
\[
\begin{aligned}
 W&=(\alpha_0,\beta_1,\alpha_1,\ldots,\beta_a,\alpha_a),
       &&a\geq1,\\
 V&=(\gamma_0,\delta_1,\gamma_1,\ldots,\delta_b,\gamma_b),
       &&b\geq1.
\end{aligned}
\]
All entries belong to \(\operatorname{E}(w)\). Define
\[
\begin{split}
 W\mathbin{\star_\eta}V\defeq
  (&\alpha_0,\beta_1,\alpha_1,\ldots,\beta_a,\,
                  \eta\land\alpha_a\land\gamma_0,\\
   &\delta_1,\gamma_1,\ldots,\delta_b,\gamma_b).
\end{split}
\]
The symbol \(\star_\eta\) is a finite word operation: it retains the
left word up to its last point, combines the three midpoint labels at
that point, and then appends the right word after its first point.
It is not an object-language connective.
The conjunction is associated as \((\eta\land\alpha_a)\land\gamma_0\).

The resulting word has \(a+b\) intervals, \(a+b+1\) point positions
and \(2(a+b)+1\) alternating point and interval positions. Its midpoint
is point position \(a\), counting from zero. Every interval label from
both inputs is retained, and both labels formerly attached to the shared
point occur in its new conjunction. In particular \(a,b\geq1\) make
the midpoint an interior position. Labels remain formulas on the whole
structure at their sole free argument; their internal quantifiers are
not restricted to either constituent interval.
\end{hbdefinition}

\begin{hblemma}{thm:p09-kamp-elementary-midpoint}{proved=annex:p9.kamp.elementary.midpoint}{The elementary midpoint concatenation}
Let \(x,y,z\) be distinct variable names, let
\(\theta\in\operatorname{E}(z)\), and let
\(\rho\in\operatorname{Ep}(x,z)\),
\(\sigma\in\operatorname{Ep}(z,y)\).
Obtain their elementary words \(W,V\) with a common label parameter
\(w\) distinct from \(x,y,z\), using
\cref{thm:p09-kamp-elementary-intervals}, and put
\(\eta\defeq\theta[w/z]\), renaming bound variables before replacing
free occurrences. Then
\[
                 \mathcal B_p(W\mathbin{\star_\eta}V;x,y)
\]
has an \(\operatorname{Ep}(x,y)\) presentation and
\[
\begin{split}
 \mathsf{Lin}\models x<y\longrightarrow
 \Bigl(&\exists z\,(x<z<y\land\theta\land\rho\land\sigma)\\
       &\longleftrightarrow
              \mathcal B_p(W\mathbin{\star_\eta}V;x,y)\Bigr).
\end{split}
\]
Every bound variable in the constituent formulas is first made fresh
from the other constituents and from \(x,y,z\). The outer
\(\exists z\) binds the original free occurrences of \(z\) in
\(\theta,\rho,\sigma\). The output has exactly \(x,y\) free; the
parameter \(w\) and the intermediate endpoint \(z\) do not remain free.
This is the elementary concatenation underlying the source's
display~(21), with both directions supplied in the annex
\cite[ch.~III, p.~79]{Kamp1968}.
Its \(\mathsf{Lin}\) scope is an additional conclusion of that
finite-realization proof.
\end{hblemma}

\begin{hbconvention}{def:p09-kamp-midpoint-families}{Finite families and their empty cases}
Let \(x,y,z\) be distinct variable names. For finite family sizes
\(m,n,p\geq0\), take
\[
\begin{gathered}
 \chi_i\in\operatorname{E}(z)\quad(1\leq i\leq m),\\
 \phi_j\text{ with property }I(x,z)\quad(1\leq j\leq n),\qquad
 \psi_k\text{ with property }I(z,y)\quad(1\leq k\leq p).
\end{gathered}
\]
The property \(I\) is exactly the guarded \(\mathsf{Com}\) property
of \cref{def:p09-kamp-interval-property}. In particular a \(\phi_j\)
has no free variables other than possibly \(x,z\), and a \(\psi_k\)
has no free variables other than possibly \(z,y\).

Write the three aggregate formulas as
\[
\begin{aligned}
 \theta(z)&\defeq
  \begin{cases}\bigwedge_{i=1}^{m}\chi_i,&m>0,\\T^K_z,&m=0,\end{cases}\\
 \Phi(x,z)&\defeq
  \begin{cases}\bigwedge_{j=1}^{n}\phi_j,&n>0,\\T^K_x,&n=0,\end{cases}\\
 \Psi(z,y)&\defeq
  \begin{cases}\bigwedge_{k=1}^{p}\psi_k,&p>0,\\T^K_y,&p=0,\end{cases}
\end{aligned}
\]
where finite nonempty conjunctions are associated from left to right.
The arguments of \(\Phi,\Psi\) display their permitted free variables;
they do not assert that every argument actually occurs freely.
The formulas \(T^K_v=P_0(v)\lor\neg P_0(v)\) are the source
tautologies already fixed in \cref{def:p09-kamp-elementary-class}.
Define the binary formula
\[
 \Gamma(x,y)\defeq
 \exists z\bigl(x<z<y\land\theta(z)\land\Phi(x,z)\land\Psi(z,y)\bigr).
\]
Bound variables inside all inputs are renamed fresh before this
construction. The outer order guard supplies both free endpoints,
including when every family is empty.

The source displays the indexed families and their three conjunctions
in (20), and reduces the midpoint family to one formula. It gives no
separate zero-family analysis there. The definitions for zero sizes
above make the handbook's empty-family reading explicit; they are
proved below together with the displayed positive-size case. A
zero-size family adds no semantic condition, but it does not delete
the existential midpoint or either strict inequality
\cite[ch.~III, p.~78]{Kamp1968}.
\end{hbconvention}

\begin{hbmethod}{def:p09-kamp-midpoint-construction}{Constructing the interval disjunction}
The input consists of the families above, their finite
\(\operatorname{E}\) constructions, and finite elementary disjunctions
witnessing each stated property \(I\). The construction uses those
witnesses; it does not decide whether a bare formula has property \(I\).

Choose \(w\) distinct from \(x,y,z\). Form \(\theta(z)\) and its renamed label
\(\eta(w)=\theta[w/z]\). Construct finite lists \(\mathcal L\) and
\(\mathcal R\) of elementary interval formulas as follows.

For \(n>0\), initialize \(\mathcal L\) with the given representation
of \(\phi_1\). For \(j=2,\ldots,n\), apply the finite fusion construction
of \cref{thm:p09-kamp-iii2-lemma4} to every pair consisting of a
current list member and a disjunct in the representation of \(\phi_j\).
Replace the list by all resulting elementary fusion disjuncts.
For \(n=0\), take the one-member list whose member is
\(\mathcal B_p((T^K_w,T^K_w,T^K_w);x,z)\).
Construct \(\mathcal R\) in the same way from
\(\psi_1,\ldots,\psi_p\) at endpoints \(z,y\); for \(p=0\), use
\(\mathcal B_p((T^K_w,T^K_w,T^K_w);z,y)\).

A supplied empty disjunction means false on the ordered endpoints.
It can first be replaced by the one false elementary pattern in
\cref{def:p09-kamp-interval-property}. Thus all lists in this method
may be taken nonempty. A singleton family needs no conjunction step.
Repeated disjuncts may be retained.

For each \(\rho\in\mathcal L\), obtain its elementary word \(W_\rho\);
for each \(\sigma\in\mathcal R\), obtain its word \(V_\sigma\).
Return the finite disjunction
\[
 \bigvee_{\rho\in\mathcal L}\ \bigvee_{\sigma\in\mathcal R}
       \mathcal B_p(W_\rho\mathbin{\star_\eta}V_\sigma;x,y).
\]
Every original choice of disjuncts on the left and on the right,
and every intervening fusion branch, is included. The final step takes
every pair from the two resulting lists. The shared midpoint fixes
the order of the two constituent pieces.
\end{hbmethod}

\begin{hblemma}{thm:p09-kamp-iii2-lemma5}{proved=annex:p9.kamp.iii2.lemma5}{Kamp's guarded interval composition lemma}
For positive \(m,n,p\), the source formula
\[
 \exists z\left(x<z<y\land
       \bigwedge_{i=1}^{m}\chi_i\land
       \bigwedge_{j=1}^{n}\phi_j\land
       \bigwedge_{k=1}^{p}\psi_k\right)
\]
has property \(I(x,y)\), under the hypotheses of
\cref{def:p09-kamp-midpoint-families}.

With the explicit empty-family convention there, the conclusion also
holds for every \(m,n,p\geq0\). More precisely, the construction in
\cref{def:p09-kamp-midpoint-construction} gives
\[
\begin{split}
 \mathsf{Com}\models x<y\longrightarrow
 \Bigl(&\Gamma(x,y)\\
 &\longleftrightarrow
   \bigvee_{\rho\in\mathcal L}\ \bigvee_{\sigma\in\mathcal R}
      \mathcal B_p(W_\rho\mathbin{\star_\eta}V_\sigma;x,y)\Bigr).
\end{split}
\]
Every disjunct has an \(\operatorname{Ep}(x,y)\) presentation.
The construction is finite. It uses the predicate letters of the
supplied elementary labels, together with the declared \(P_0\)
tautology or falsehood when an empty-family or empty-disjunction
convention requires it. Fresh letters used to exhibit the source
substitution construction are eliminated.

The given \(I\)-representations and their conjunctions retain
\(\mathsf{Com}\) scope. The elementary midpoint identity is applied
to their actual elementary disjuncts and then restricted to complete
orders. This proves the source's conclusion without promoting an
arbitrary \(I\)-representation to \(\mathsf{Lin}\)
\cite[ch.~III, lemma~5, pp.~78--79]{Kamp1968}.
\end{hblemma}

\begin{hbexample}{def:p09-kamp-midpoint-one-piece}{Two one-interval patterns and their shared point}
Let \(A,B,C,D,E,F\) be unary predicate letters and let
\(H\in\operatorname{E}(w)\). View \(A(w),\ldots,F(w)\) as elementary
labels. The two words \((A,B,C)\) and \((D,E,F)\), joined with
midpoint label \(H\), give
\[
                         (A,B,H\land C\land D,E,F).
\]
With fresh bound variables this pattern is
\[
\begin{split}
 A(x)\land\exists z\bigl(&x<z<y\land
          \forall u\,(x<u<z\to B(u))\\
      &{}\land H(z)\land C(z)\land D(z)\\
      &{}\land\forall v\,(z<v<y\to E(v))\land F(y)\bigr).
\end{split}
\]
All three middle labels concern the same chosen \(z\).
The letter \(E\) in this example is a unary predicate; the class
\(\operatorname{E}\) continues to denote the class of formulas.

On \(0<1<2\), take \(A,B,C,D,E,F\) true everywhere. The formula at
\((0,2)\) holds exactly when \(H\) holds at 1; its two open intervals
are empty. Requiring \(H\) only at either outer endpoint would change
the condition. If all three input families are empty, the conclusion
instead represents precisely \(\exists z\,(x<z<y)\). It is false
at adjacent endpoints and true at \((0,2)\) in this three-point order.
An empty family therefore leaves a substantive strict-witness condition.
\end{hbexample}

\paragraph{The next source boundary.}
Lemma~6 begins on printed p.~79. It fixes one basic pattern and four
strictly increasing points, assumes realizations from the first point
to the second and fourth, and characterizes realization to the third
using a guarded past closure. That separate finite-pattern argument
is the next dependency. Lemma~5 supplies the midpoint composition
above; it does not complete the remaining interval lemmas or
Theorem~III.2 \cite[ch.~III, p.~79]{Kamp1968}.
